\begin{frame}{하나의 시뮬레이션을 따라가보자}
  루트(돌 5개, \textbf{내 차례})에서 시작:
  \begin{enumerate}
    \item \kb{선택} --- 루트에 자식이 \textbf{하나도 없으므로}
      루트에서 즉시 정지
    \item \kb{확장} --- 아직 안 가본 수(예: ``2개 가져가기'')를
      무작위로 하나 골라 자식으로 추가
    \item \kb{시뮬레이션} --- 새 노드(돌 3개, \textbf{상대 차례})부터
      무작위 플레이로 게임 끝까지 두어 승자를 판정
    \item \kb{역전파} --- 루트와 그 새 자식, \textbf{두 노드 모두}에
      ``방문 1회''를 더하고, 승자가 ``그 노드로 \textbf{두어
      도달한 플레이어}''라면 승 1회도 더함
  \end{enumerate}
  \vskip0.6em
  \begin{block}{``모든 수를 최소 한 번'' 보장은 여기서 온다}
    \textbf{첫 3회} 시뮬레이션은 세 갈림의 각각을 한 번씩
    ``열어보는'' 데 쓰임 --- \textbf{방문 0인 자식은 항상
    우선}되기 때문. 4회차부터 선택 단계의 UCT가 실제로
    \textbf{작용}하기 시작한다.
  \end{block}
\end{frame}
